\textbf{Limitations of the evidence.}
\begin{itemize}
\item \emph{Small samples.} $n\le5$ per SWE cell and 3 seeds per market cell. The design detects only large
effects (about 0.6 absolute), and none of the SWE contrasts reached one. Market results are directions,
not estimates.
\item \emph{One task, one agent model.} All SWE results come from task T3 and all untrusted agents are
\texttt{gemma4:31b}. Low competence (6/49 correct) dominated the outcomes; a stronger model or easier task
may change every SWE conclusion.
\item \emph{Induced, not learned, collusion.} Colluders were told their goal and channel. Collusion that
emerges under optimisation pressure may be quieter, and less likely to be announced in chat.
\item \emph{Stochastic pairing.} LLM sellers are not deterministic, so paired market runs diverge before the
intervention. We used a before/after difference against the pair, which only partly corrects for this.
\item \emph{Late, short interventions.} The market monitor acted from round 10, so interventions ran for
1--5 of 15 rounds.
\item \emph{Deviations.} One excluded SWE episode was not replaced as pre-registered, and the \Mtwo{} LLM
judge was not replayed. Market results (RQ4) are exploratory in the pre-registration's terms.
\end{itemize}

\textbf{Constraints of the setup.}
\begin{itemize}
\item \emph{Structural trust assumptions.} Agents cannot write to the evidence log, reach the monitor, or
bypass the release or order gate. These hold in our harness by construction; a deployed system has to
enforce them.
\item \emph{The reviewer was honest by construction} and became the strongest safeguard. A coalition that
includes the reviewer is outside this study.
\item \emph{Verification needs observable outcomes.} Hidden tests must cover the defect (our \Mthree{} sees
half of them), and the market check needs trades: a cartel that refuses to sell produced no prices to
verify.
\item \emph{Free endpoints.} Rate limits and latency forced one pinned model and small runs; budget, not
design, set the sample size.
\item \emph{Same-family trusted roles.} The \Mtwo{} judge and \qa{} wording share one model; the trusted
side's verdicts were deterministic, so the capability gap between trusted and untrusted components was not
varied.
\end{itemize}

\textbf{Open questions.}
\begin{itemize}
\item Does a trusted verifier with a \textsc{Hold} right, rather than a voice, reduce harm, and at what
utility cost?
\item Would a stronger or goal-directed colluder accept a stranger's offer? Is the null probe evidence of
honesty, policy-following, or simply ignoring unsolicited messages?
\item How much entry is needed to make incumbents compete, rather than only giving buyers an outside option?
\item When the inner monitor is itself a capable, untrusted model that agents can see and persuade, can an
outer layer of weaker but trusted checks detect that the monitor has joined the coalition?
\item How should harm verification and collusion detection be combined when honest agents also fail?
\end{itemize}
